\section{Vị trí tương đối giữa hai đường thẳng, góc và khoảng cách}

\subsection{Đề 1}
\Opensolutionfile{ans}[ans/ans-0-C7B10-De1-LC]
\caulc
%cau 1
\begin{ex}%[0H9H3-3]
Xét vị trí tương đối của hai đường thẳng $d_1 \colon x-2 y+1=0$ và $d_2 \colon -3 x+6 y-10=0$.
\choice
{Trùng nhau}
{\True Song song}
{Vuông góc với nhau}
{Cắt nhau nhưng không vuông góc nhau}
\loigiai{
Ta thấy $ \dfrac{1}{-3}=\dfrac{-2}{6} \neq \dfrac{1}{-10} $ nên $ d_1 \parallel d_2 $.
}
\end{ex}

%cau 2
\begin{ex}%[0H9H3-3]
Xét vị trí tương đối của hai đường thẳng $d_1 \colon \dfrac{x}{3}-\dfrac{y}{4}=1$ và $d_2 \colon 3 x+4 y-10=0$.
\choice
{Trùng nhau}
{Song song}
{\True Vuông góc với nhau}
{Cắt nhau nhưng không vuông góc nhau}
\loigiai{
$ d_1 $ có một vectơ pháp tuyến $ \overrightarrow{n}_1=\left(\dfrac{1}{3} ;-\dfrac{1}{4}\right)  $.
\\
$ d_2 $ có một vectơ pháp tuyến $ \overrightarrow{n}_2=(3 ; 4)  $.
\\
Ta thấy
$$
	\overrightarrow{n}_1 \cdot \overrightarrow{n}_2 = 1-1 = 0
$$
nên hai đường này vuông góc nhau.
}
\end{ex}

%cau 3
\begin{ex}%[0H9H3-3]
Cho hai đường thẳng $\left(d_1\right) \colon m x+y=m+1$, $\left(d_2\right) \colon x+m y=2$ cắt nhau khi và chỉ khi
\choice
{$m \neq 2$}
{$m \neq \pm 1$}
{\True $m \neq 1$}
{$m \neq-1$}
\loigiai{
Ta có
$$
	\left(d_1\right) \cap \left(d_2\right)
	\Leftrightarrow
	\heva{&m x+y=m+1 &(1) \\ &x+m y=2 &(2)} \quad \text{có một nghiệm.}
$$ 
Thay $(2)$ vào $(1) \Rightarrow m(2-m y)+y=m+1 \Leftrightarrow\left(1-m^2\right) y=1-m(*)$.
\\
Hệ phương trình có một nghiệm $\Leftrightarrow(*)$ có một nghiệm 
$\Leftrightarrow
\heva{&1-m^2 \neq 0 \\ &m-1 \neq 0}
\Leftrightarrow m \neq 1$.
}
\end{ex}

%cau 4
\begin{ex}%[0H9H3-3]
Giao điểm $M$ của $(d) \colon \heva{&x=1-2 t \\ &y=-3+5 t}$ và $\left(d'\right) \colon 3 x-2 y-1=0$ là
\choice
{$M\left(2 ;-\dfrac{11}{2}\right)$}
{$M\left(0 ; \dfrac{1}{2}\right)$}
{\True $M\left(0 ;-\dfrac{1}{2}\right)$}
{$M\left(-\dfrac{1}{2} ; 0\right)$}
\loigiai{
Ta có
$$
	(d) \colon \heva{&x=1-2 t \\ &y=-3+5 t}
	\Rightarrow
	(d) \colon 5 x+2 y+1=0.
$$
Do $ M $ là giao điểm của $ (d) $ và $ (d') $ nên tọa độ điểm $ M $ là nghiệm của hệ phương trình
$$
	\heva{&3 x-2 y-1=0 \\ &5 x+2 y+1=0}
	\Leftrightarrow
	\heva{&x=0 \\ &y=-\dfrac{1}{2}.}
$$
Vậy $M\left(0 ;-\dfrac{1}{2}\right)$.
}
\end{ex}

%cau 5
\begin{ex}%[0H9H3-3]
Hai đường thẳng $\left(d_1\right) \colon \heva{&x=-2+5 t \\ &y=2 t}$ và $\left(d_2\right) \colon 4 x+3 y-18=0$ cắt nhau tại điểm có tọa độ
\choice
{\True $(2 ; 3)$}
{$(3 ; 2)$}
{$(1 ; 2)$}
{$(2 ; 1)$}
\loigiai{
Ta có
$$
	\left(d_1\right) \colon \heva{&x=-2+5 t \\&y=2 t}
	\Rightarrow
	\left(d_1\right) \colon 2 x-5 y+4=0
$$
Gọi $ M $ là giao điểm của $ (d_1) $ và $ (d_2) $ nên tọa độ $ M $ là nghiệm của hệ phương trình
$$
	\heva{&2 x-5 y+4=0 \\ &4 x+3 y-18=0}
	\Leftrightarrow
	\heva{&x=2 \\ &y=3.}
$$
Vậy $ M(2;3) $.
}
\end{ex}

%cau 6
\begin{ex}%[0H9H3-3]
Cho 4 điểm $A(1 ; 2)$, $B(4 ; 0)$, $C(1 ;-3)$, $D(7 ;-7)$. Xác định vị trí tương đối của hai đường thẳng $A B$ và $C D$.
\choice
{\True Song song}
{Cắt nhau nhưng không vuông góc}
{Trùng nhau}
{Vuông góc nhau}
\loigiai{
Ta có $\overrightarrow{A B}=(3 ;-2)$, $\overrightarrow{C D}=(6 ;-4)$.
\\
Mà $\dfrac{3}{6}=\dfrac{-2}{-4}$ nên $ AB \parallel CD $.
}
\end{ex}

%cau 7
\begin{ex}%[0H9V3-3]
Trong mặt phẳng với hệ tọa độ $O x y$, cho ba đường thẳng lần lượt có phương trình $d_1 \colon 3 x-4 y+15=0$, $d_2 \colon 5 x+2 y-1=0$ và $d_3 \colon m x-(2 m-1) y+9 m-13=0$. Tìm tất cả các giá trị của tham số $m$ để ba đường thẳng đã cho cùng đi qua một điểm.
\choice
{$m=\dfrac{1}{5}$}
{$m=-5$}
{$m=-\dfrac{1}{5}$}
{\True $m=5$}
\loigiai{
Xét hệ phương trình tọa độ giao điểm của $ (d_1) $ và $ (d_2) $
$$
	\heva{&3 x-4 y+15=0 \\&5 x+2 y-1=0}
	\Leftrightarrow
	\heva{&x=-1 \\&y=3.}
$$
Suy ra giao điểm của $ (d_1) $ và $ (d_2) $ là $ A(-1 ; 3) $.
Để $ (d_1) $, $ (d_2) $ và $ (d_3) $ cùng đi qua một điểm thì
$$
	A \in (d_3)
	\Leftrightarrow
	-m-6 m+3+9 m-13=0
	\Leftrightarrow
	m=5.
$$
Vậy $ m=5 $ thỏa yêu cầu bài toán.
}
\end{ex}

%cau 8
\begin{ex}%[0H9V3-3]
Với giá trị nào của $m$ thì ba đường thẳng $d_1 \colon 3 x-4 y+15=0$, $d_2 \colon 5 x+2 y-1=0$ và $d_3 \colon m x-4 y+15=0$ đồng quy?
\choice
{$ m=-5 $}
{$ m=5 $}
{\True $ m=3 $}
{$ m=-3 $}
\loigiai{
Xét hệ phương trình tọa độ giao điểm của $ (d_1) $ và $ (d_2) $
$$
	\heva{&3 x - 4 y + 1 5 = 0 \\&5 x+2 y-1=0}
	\Leftrightarrow
	\heva{&x=-1 \\&y=3.}
$$
Suy ra giao điểm của $ (d_1) $ và $ (d_2) $ là $ A(-1 ; 3) $.
Để $ (d_1) $, $ (d_2) $ và $ (d_3) $ đồng quy thì
$$
	A \in (d_3)
	\Leftrightarrow
	-m-12+15=0
	\Leftrightarrow
	m=3
$$
Vậy $ m=3 $ thỏa yêu cầu bài toán.
}
\end{ex}

%cau 9
\begin{ex}%[0H9H3-3]
Cho 4 điểm $A(-3 ; 1)$, $B(-9 ;-3)$, $C(-6 ; 0)$, $D(-2 ; 4)$. Tìm tọa độ giao điểm của 2 đường thẳng $A B$ và $C D$.
\choice
{$(-6 ;-1)$}
{\True $(-9 ;-3)$}
{$(-9 ; 3)$}
{$(0 ; 4)$}
\loigiai{
Do $\overrightarrow{A B}=(-6 ;-4)$ nên đường thẳng 
AB có một vectơ pháp tuyến $\overrightarrow{n}_{A B}=(2 ;-3)$.
\\
Suy ra phương trình đường thẳng $ AB $ là $ 2 x-3 y=-9$.
\\
Tương tự, ta cũng có phương trình đường thẳng $ CD $ là $ x-y=-6 $.
\\
Tọa độ giao điểm $N$ của $ AB $ và $ CD $ là nghiệm của hệ phương trình
$$
	\heva{&2 x-3 y=-9 \\ &x-y=-6}
	\Leftrightarrow
	\heva{&x=-9 \\ &y=-3}
$$
Vậy $ N(-9 ;-3) $.
}
\end{ex}

%cau 10
\begin{ex}%[0H9V3-3]
Hai đường thẳng $d_1 \colon m x+y=m-5$, $d_2 \colon x+m y=9$ cắt nhau khi và chỉ khi
\choice
{$m \neq-1$}
{$m \neq 1$}
{\True $m \neq \pm 1$}
{$m \neq 2$}
\loigiai{
$d_1$ và $d_2$ theo thứ tự nhận các vectơ $\overrightarrow{n}_1=(m ; 1)$, $\overrightarrow{n}_2=(1 ; m)$ làm vectơ pháp tuyến.
\\
$d_1$ và $d_2$ cắt nhau $\Leftrightarrow \overrightarrow{n}_1$ và $\overrightarrow{n}_2$ không cùng phương $\Leftrightarrow m \cdot m \neq 1 \cdot 1 \Leftrightarrow m \neq \pm 1$.
}
\end{ex}

%cau 11
\begin{ex}%[0H9H3-3]
Trong mặt phẳng với hệ tọa độ $O x y$, cho hai đường thẳng có phương trình $d_1 \colon m x+(m-1) y+2 m=0$ và $d_2 \colon 2 x+y-1=0$. Nếu $d_1$ song song $d_2$ thì
\choice
{\True $m=2$}
{$m=-1$}
{$m=-2$}
{$m=1$}
\loigiai{
Để $d_1 \colon m x+(m-1) y+2 m=0$ và $d_2 \colon 2 x+y-1=0$ song song với nhau thì
$$
	\dfrac{m}{2}=\dfrac{m-1}{1} \neq \dfrac{2 m}{-1}
	\Leftrightarrow
	\heva{&-1 \neq 2 \\&m=2 m-2}
	\Leftrightarrow
	m=2.
$$
}
\end{ex}

%cau 12
\begin{ex}%[0H9H3-3]
Tính góc giữa hai đường thẳng $\Delta \colon x-\sqrt{3} y+2=0$ và $\Delta' \colon x+\sqrt{3} y-1=0$.
\choice
{$90^{\circ}$}
{$120^{\circ}$}
{\True $60^{\circ}$}
{$30^{\circ}$}
\loigiai{
Đường thẳng $\Delta$ có vectơ pháp tuyến $\overrightarrow{n}=(1 ;-\sqrt{3})$.
\\
Đường thẳng $\Delta'$ có vectơ pháp tuyến $\overrightarrow{n'}=(1 ; \sqrt{3})$.
\\
Gọi $\alpha$ là góc giữa hai đường thẳng $\Delta$, $\Delta' $. Ta có
$$
	\cos \alpha
	=\left|\cos \left(\overrightarrow{n}, \overrightarrow{n}'\right)\right|
	=\dfrac{|1-3|}{\sqrt{1+3} \cdot \sqrt{1+3}}=\dfrac{1}{2}
	\Rightarrow \alpha=60^{\circ}
$$
}
\end{ex}

\Closesolutionfile{ans}
\indapan{6}{ans/ans-0-C7B10-De1-LC}

\cauds
\Opensolutionfile{ans}[ans/ans-0-C7B10-De1-DS]
\begin{ex}%[0H9H3-3]
Các mệnh đề sau đúng hay sai?
\choiceTF
{\True $ d_1 \colon x+4 y-1=0 $ cắt $ d_2 \colon 2 x-3 y+5=0 $}
{\True $ m_1 \colon \heva{&x=5+3 t \\&y=-5-4 t} $ song song $ m_2 \colon 8 x+6 y+1=0 $}
{\True $ a_1 \colon \heva{&x=1+t \\&y=-3+3 t} $ trùng $ a_2 \colon \heva{&x=2+2 k \\&y=6 k} $ (với $ t $, $ k $ là các tham số).}
{$ \Delta_1 \colon x+y-1=0 $ và $ \Delta_2 \colon x+2=0 $; góc giữa hai đường thẳng là $ 30^{\circ} $}
\loigiai{
Ta thấy
\begin{itemchoice}
	\itemch Hai đường thẳng $ d_1 \colon x+4 y-1=0 $ và $ d_2 \colon 2 x-3 y+5=0 $ cắt nhau vì $ \dfrac{1}{2} \ne \dfrac{4}{-3} $.
	\itemch $ m_1 \colon \heva{&x=5+3 t \\&y=-5-4 t} $ song song $ m_2 \colon 8 x+6 y+1=0 $ vì vectơ chỉ phương $ \overrightarrow{u}_1 = (3;-4) $ của $ m_1 $ vuông góc với vectơ pháp tuyến $ \overrightarrow{u}_2 = (8;6) $ của $ m_2 $ ($ 3 \cdot 8 + (-4) \cdot 6 = 0 $).
	\itemch Ta có
	\begin{align*}
		&a_1 \colon \heva{&x=1+t \\&y=-3+3 t}
		\Rightarrow
		3x-y-6=0; \\
		&a_2 \colon \heva{&x=2+2 k \\&y=6 k}
		\Rightarrow
		3x-y-6=0.
	\end{align*}
	$ a_1 $ và $ a_2 $ có phương trình tổng quát giống nhau nên chúng trùng nhau.
	\itemch $ \Delta_1 \colon x+y-1=0 $ có vectơ pháp tuyến là $ \overrightarrow{n}_1=(1;1) $.
	\\
	$ \Delta_2 \colon x+2=0 $ có vectơ pháp tuyến là $ \overrightarrow{n}_2=(1;0) $.
	\\
	Gọi $ \alpha $ là góc giữa hai đường $ \Delta_1 $ và $ \Delta_2 $, ta có
	$$
		\cos \alpha
		=\left|\cos \left(\overrightarrow{n}_1, \overrightarrow{n}_2\right)\right|
		=\dfrac{|1+0|}{\sqrt{1+1} \cdot \sqrt{1+0}}=\dfrac{1}{\sqrt{2}}
		\Rightarrow
		\alpha=45^{\circ}
	$$
\end{itemchoice}
}
\end{ex}

%cau 2
\begin{ex}%[0H9H3-3]
Cho hai đường thẳng $ \Delta_1 \colon \heva{&x=2+5 t \\&y=3-6 t} $ và $ \Delta_2 \colon \heva{&x=7+5 t' \\&y=-3+6 t'} $.
\\
Các mệnh đề sau đúng hay sai?
\choiceTF
{\True Hai đường thẳng $\Delta_1$, $\Delta_2$ lần lượt có vectơ chỉ phương $\overrightarrow{u}_1=(5 ;-6)$, $\overrightarrow{u}_2=(5 ; 6)$}
{Hai đường thẳng $\Delta_1$, $\Delta_2$ song song}
{$M(7 ; 3)$ là tọa độ giao điểm hai đường $\Delta_1$, $\Delta_2$}
{$\Delta_1$, $\Delta_2$ vuông góc với nhau}
\loigiai{
\begin{itemchoice}
	\itemch Hai đường thẳng $\Delta_1$, $\Delta_2$ lần lượt có vectơ chỉ phương $\overrightarrow{u}_1=(5 ;-6)$, $\overrightarrow{u}_2=(5 ; 6)$.
	
	\itemch Xét $\overrightarrow{u}_1=(5 ;-6)$, $\overrightarrow{u}_2=(5 ; 6)$ với $5 \cdot 6 \neq -6 \cdot 5$ nên hai vectơ này không cùng phương.
	\\
	Vì vậy hai đường thẳng $\Delta_1$, $\Delta_2$ cắt nhau.
	
	\itemch Giải hệ phương trình
	$$
		\heva{&2+5 t=7+5 t' \\ &3-6 t=-3+6 t'}
		\Leftrightarrow
		\heva{&5 t-5 t'=5 \\ &-6 t-6 t'=-6}
		\Leftrightarrow
		\heva{&t=1 \\ &t'=0}
		\Rightarrow
		M(7 ;-3).
	$$
	$ M(7 ;-3) $ là tọa độ giao điểm hai đường $\Delta_1$, $\Delta_2$.
	
	\itemch Ta có
	$$
		\overrightarrow{u}_1 \cdot \overrightarrow{u}_2=5 \cdot 5-6 \cdot 6=-11 \neq 0.
	$$
	Suy ra hai đường thẳng đã cho chỉ cắt nhau mà không vuông góc.
\end{itemchoice}
}
\end{ex}

\begin{ex}%[0H9H3-3]
Các mệnh đề sau đúng hay sai?
\choiceTF
{$M(2 ;-1)$; $ \Delta \colon 3 x-4 y-12=0$, khi đó $\mathrm{d}(M, \Delta)=\dfrac{3}{5}$}
{\True $M(4 ;-5)$; $ \Delta \colon \heva{&x=2 t \\ &y=2+3 t}$, khi đó $\mathrm{d}(M, \Delta)=2 \sqrt{13}$}
{\True $\Delta_1 \colon 7 x+y-3=0$ và $\Delta_2 \colon 7 x+y+12=0$ song song với nhau.}
{$\Delta_1 \colon 7 x+y-3=0$ và $\Delta_2 \colon 7 x+y+12=0$ khi đó $d\left(\Delta_1, \Delta_2\right)=\dfrac{\sqrt{2}}{2}$}
\loigiai{
\begin{itemchoice}
	\itemch Ta có
	$$
		\mathrm{d}(M, \Delta)=\dfrac{|3 \cdot 2-4 \cdot(-1)-12|}{\sqrt{3^2+(-4)^2}}
		=\dfrac{2}{5}.
	$$
	
	\itemch Ta có
	$$
		\Delta\colon\heva{&x=2 t \\ &y=2+3 t} 
		\Rightarrow 
		\Delta\colon\dfrac{x}{2}=\dfrac{y-2}{3}
		\Rightarrow \Delta \colon 3 x-2 y+4=0.
	$$
	Do đó
	$$
		\mathrm{d}(M, \Delta)=\dfrac{|3 \cdot 4-2 \cdot(-5)+4|}{\sqrt{3^2+(-2)^2}}
		=2 \sqrt{13}.
	$$
	
	\itemch $\Delta_1 \colon 7 x+y-3=0$ và $\Delta_2 \colon 7 x+y+12=0$ song song vì $ \dfrac{7}{7}=\dfrac{1}{1} \ne \dfrac{-3}{12} $.
	
	\itemch Lấy $M(0 ; 3) \in \Delta_1$, khi đó
	$$
		\mathrm{d}\left(\Delta_1, \Delta_2\right)
		=\mathrm{d}\left(M, \Delta_2\right)
		=\dfrac{|7.0+3+12|}{\sqrt{7^2+1^2}}
		=\dfrac{3 \sqrt{2}}{2}.
	$$
\end{itemchoice}
}
\end{ex}

\begin{ex}%[0H9H3-3]
Cho $\Delta_1 \colon \heva{&x=3-t \\ &y=2-t}$, $\Delta_2 \colon \heva{&x=1+2 t' \\ &y=1-3 t'}$.
\\
Các mệnh đề sau đúng hay sai?
\choiceTF
{\True $\Delta_1$ có vectơ chỉ phương $\overrightarrow{u_1}=(-1 ;-1)$}
{\True $\Delta_2$ có vectơ chỉ phương $\overrightarrow{u_2}=(2 ;-3)$}
{Hai đường thẳng $\Delta_1$, $\Delta_2$ song song}
{$\Delta_1$, $\Delta_2$ cắt nhau tại điểm có tọa độ $\left(\dfrac{7}{3} ; \dfrac{2}{3}\right)$}
\loigiai{
\begin{itemchoice}
	\itemch $\Delta_1 \colon \heva{&x=3-t \\ &y=2-t}$ có vectơ chỉ phương $\overrightarrow{u_1}=(-1 ;-1)$.
	
	\itemch $\Delta_2 \colon \heva{&x=1+2 t' \\ &y=1-3 t'}$ có vectơ chỉ phương $\overrightarrow{u_2}=(2 ;-3)$.
	
	\itemch Xét $\overrightarrow{u}_1=(-1 ;-1)$ và $\overrightarrow{u}_2=(2 ;-3)$ với $-1 \cdot (-3) \neq-1 \cdot 2$ nên hai vectơ này không cùng phương.
	\\
	Do đó hai đường $\Delta_1$, $\Delta_2$ cắt nhau.
	
	\itemch Xét hệ hai phương trình
	$$
		\heva{&3-t=1+2 t' \\ &2-t=1-3 t'}
		\Leftrightarrow
		\heva{&-t-2 t'=-2 \\ &-t+3 t'=-1}
		\Leftrightarrow
		\heva{&t=\dfrac{8}{5} \\ &t'=\dfrac{1}{5}}
		\Rightarrow
		\heva{&x=\dfrac{7}{5} \\ &y=\dfrac{2}{5}.}
	$$
	Vậy $\Delta_1$, $\Delta_2$ cắt nhau tại điểm có tọa độ $\left(\dfrac{7}{5} ; \dfrac{2}{5}\right)$.
\end{itemchoice}
}
\end{ex}

\Closesolutionfile{ans}
\indapan{2}{ans/ans-0-C7B10-De1-DS}


\Opensolutionfile{ans}[ans/ans-0-C7B10-De1-KQ]
\caukq
%cau 1
\begin{ex}%[0H9H3-3]
Cho đường thẳng $\Delta \colon 3 x+4 y-6=0$ và $\Delta' \colon x+y=1$. Điểm $M$ thuộc $\Delta'$ sao cho khoảng cách từ $M$ đến $\Delta$ bằng $\dfrac{4}{5}$. Có bao nhiêu điểm $ M $ thỏa mãn yêu cầu bài toán.
\shortans{$ 2 $}
\loigiai{
Viết phương trình tham số
$$
	\Delta' \colon \heva{&x=t \\ &y=1-t}
$$
Gọi $M(t ; 1-t) \in \Delta'$. Khi đó
$$
	d(M, \Delta)
	=\dfrac{|3 t+4(1-t)-6|}{\sqrt{3^2+4^2}}
	=\dfrac{|-t-2|}{5}=\dfrac{4}{5}
	\Rightarrow
	|t+2|=4.
	\Leftrightarrow
	\hoac{&t=2 \\&t=-6}
$$
Phương trình có $ 2 $ nghiệm nên ta tìm được $ 2 $ điểm thỏa mãn đề bài.
}
\end{ex}

%cau 2
\begin{ex}%[0H9V3-3]
Cho tam giác $A B C$ có phương trình đường thẳng chứa các cạnh $A B$, $A C$, $B C$ lần lượt là $x+2 y-1=0$; $x+y+2=0$; $2 x+3 y-5=0$. Tính diện tích tam giác $A B C$.
\shortans{$ 18 $}
\loigiai{
Tọa độ của điểm $A$ là nghiệm của hệ phương trình
$$
	\heva{&x+2 y-1=0 \\ &x+y+2=0}
	\Leftrightarrow
	\heva{&x=-5 \\ &y=3}
$$
Suy ra điểm $A$ có tọa độ là $(-5 ; 3)$.
\\
Gọi $A H$ là đường cao kẻ từ $A$ của tam giác $A B C$ $(H \in B C)$. Ta có
$$
	A H=d(A, B C)
	=\dfrac{|2 \cdot(-5)+3 \cdot 3-5|}{\sqrt{2^2+3^2}}
	=\dfrac{6 \sqrt{13}}{13} .
$$
Từ các phương trình đường thẳng chứa các cạnh của tam giác $A B C$ ta tính đuợc tọa độ của điểm $B$ và điểm $C$ lần lượt là $(7 ;-3)$, $(-11 ; 9)$.
\\
Do đó, độ dài đoạn thẳng $B C$ là $6 \sqrt{13}$.
\\
Diện tích tam giác $ ABC $ bằng
$$
	\dfrac{1}{2} \cdot \dfrac{6 \sqrt{13}}{13} \cdot 6 \sqrt{13}=18.
$$
}
\end{ex}

%cau 3
\begin{ex}%[0H9H3-3]
Tìm $m$ để hai đường thẳng sau vuông góc với nhau $\Delta_1 \colon x-m y+1=0$; $\Delta_2 \colon 2 x+3 y+m=0$. Làm tròn kết quả đến hàng phần trăm.
\shortans{$ 0{,}67 $}
\loigiai{
Vectơ pháp tuyến của đường thẳng $\Delta_1 \colon x-m y+1=0$ và đường thẳng $\Delta_2 \colon 2 x+3 y+m=0$ lần lượt là $\overrightarrow{n_1}=(1 ;-m)$, $\overrightarrow{n_2}=(2 ; 3)$.
\\
Để đường thẳng $\Delta_1$ và $\Delta_2$ vuông góc với nhau thì
$$
	\overrightarrow{n_1} \perp \overrightarrow{n_2}
	\Leftrightarrow 
	\overrightarrow{n_1} \cdot \overrightarrow{n_2}=0
	\Leftrightarrow
	1 \cdot 2-m \cdot 3=0
	\Leftrightarrow
	m=\dfrac{2}{3} \approx 0{,}67.
$$
}
\end{ex}

%cau 4
\begin{ex}%[0H9V3-3]
Cho các đường thẳng $d_1 \colon x+y+3=0$, $d_2 \colon x-y-4=0$ và $d_3 \colon x-2 y=0$. Gọi điểm $M$ trên $d_3$ sao cho khoảng cách từ $M$ đến $d_1$ bằng hai lần khoảng cách từ $M$ đến $d_2$. Có bao nhiêu điểm $ M $ thỏa mãn yêu cầu này?
\shortans{$ 2 $}
\loigiai{
Ta có điểm $M$ thuộc đường thẳng $d_3$ khi và chỉ khi $M(2 t ; t)$ với $t$ là tham số.
\\
Khoảng cách từ $M$ tới $d_1$ bằng hai lần khoảng cách từ $M$ tới $d_2$ nên
$$
	\dfrac{|2 t+t+3|}{\sqrt{1^2+1^2}}
	=2 \cdot \dfrac{|2 t-t-4|}{\sqrt{1^2+(-1)^2}}
	\Leftrightarrow |3 t+3|=|2 t-8| 
	\Leftrightarrow 
	\hoac{&t=1 \\&t=-11}
$$
Vậy $M(2 ; 1)$ hoặc $M(-22 ;-11)$. Ta tìm được $ 2 $ điểm $ M $ thỏa yêu cầu bài toán.
}
\end{ex}

%cau 5
\begin{ex}%[0H9H3-3]
Tìm tham số $m$ để hai đường thẳng
$\Delta_1 \colon 2 x+\left(m^2+1\right) y-3=0$ và
$\Delta_2 \colon x+m y-100=0$
song song.
\shortans{$ 1 $}
\loigiai{
Để hai đường $ \Delta_1 $ và $ \Delta_2 $ song song thì
$$
	\dfrac{1}{2} = \dfrac{m}{m^2+1} \ne \dfrac{-100}{-3} \quad (1)
$$
Giải phương trình
$$
	\dfrac{1}{2} = \dfrac{m}{m^2+1}
	\Leftrightarrow
	m^2-2m+1 = 0
	\Leftrightarrow
	m=1.
$$
Thay $ m=1 $ vào biểu thức $ (1) $, ta thấy thỏa mãn.
}
\end{ex}

%cau 6
\begin{ex}%[0H9H3-3]
Định $m$ để hai đường thẳng $\Delta_1 \colon 2 x-3 y+4=0$ và $\Delta_2 \colon \heva{&x=2-3 t \\ &y=1-4 m t}$ vuông góc với nhau. Làm tròn kết quả đến $ 1 $ chữ số thập phân.
\shortans{$ -1{,}1 $}
\loigiai{
$\Delta_1$, $\Delta_2$ có hai vectơ pháp tuyến là $\overrightarrow{n}_1=(2 ;-3)$, $\overrightarrow{n}_2=(4 m ;-3)$.
\\
Ta có
$$
	\Delta_1 \perp \Delta_2 
	\Leftrightarrow 
	\overrightarrow{n}_1 \cdot \overrightarrow{n}_2=0
	\Leftrightarrow 
	2 \cdot 4 m+(-3) \cdot(-3)=0
	\Leftrightarrow m=-\frac{9}{8} \approx -1{,}1.
$$
}
\end{ex}

\Closesolutionfile{ans}
\indapan{6}{ans/ans-0-C7B10-De1-KQ}
\newpage
\subsection{Đề 2}
\Opensolutionfile{ans}[ans/ans-0-C7B10-De2-LC]
\caulc
%cau 1
\begin{ex}%[0H9H3-3]
Xét vị trí tương đối của hai đường thẳng $d_1 \colon 3 x-2 y-6=0$ và $d_2 \colon 6 x-2 y-8=0$.
\choice
{Trùng nhau}
{Song song}
{Vuông góc với nhau}
{\True Cắt nhau nhưng không vuông góc nhau}
\loigiai{
Đường thẳng $ d_1 \colon 3 x-2 y-6=0 $ có một vectơ pháp tuyến $ \overrightarrow{n}_1=(3 ;-2) $.
\\
Đường thẳng $ d_2 \colon 6 x-2 y-8=0 $ có một vectơ pháp tuyến $ \overrightarrow{n}_2=(6 ;-2) $.
\\
Ta thấy
$$
	\heva{&\dfrac{3}{6} \neq \dfrac{-2}{-2}
	\\&\overrightarrow{n}_1 \cdot \overrightarrow{n}_2 \neq 0.}
$$
Suy ra $ d_1 $ và $ d_2 $ cắt nhau nhưng không vuông góc.
}
\end{ex}

%cau 2
\begin{ex}%[0H9H3-3]
Xét vị trí tương đối của hai đường thẳng $d_1 \colon \heva{&x=-3+4 t \\ &y=2-6 t}$ và $d_2 \colon \heva{&x=2-2 t' \\ &y=-8+4 t'}$.
\choice
{Trùng nhau}
{\True Song song}
{Vuông góc với nhau}
{Cắt nhau nhưng không vuông góc nhau}
\loigiai{
Đường thẳng $d_1 \colon \heva{&x=-3+4 t \\ &y=2-6 t}$ có một vectơ chỉ phương $ \overrightarrow{u}_1=(2 ;-3) $ và đi qua điểm $ A(-3 ; 2) $.
\\
Đường thẳng $d_2 \colon \heva{&x=1-2 t' \\&y=4+3 t'}$ có một vectơ chỉ phương $ \overrightarrow{u}_2=(-2 ;3) $.
\\
Ta thấy
$$
	\heva{&\dfrac{2}{-2}=\dfrac{-3}{3} \\
	&A \notin d_2.}
$$
Suy ra hai đường $ d_1 $ và $ d_2 $ song song với nhau.
}
\end{ex}

%cau 3
\begin{ex}%[0H9H3-3]
Phương trình nào sau đây biểu diển đường thẳng \textbf{không} song song với đường thẳng $(d): y=2 x-1$?
\choice
{$2 x-y+5=0$}
{$2 x-y-5=0$}
{$-2 x+y=0$}
{\True $2 x+y-5=0$}
\loigiai{
Ta có $(d)\colon y=2 x-1 \Rightarrow(d)\colon 2 x-y-1=0$.
Dễ thấy $ (d) $ không song song với đường thẳng có phương trình $2 x+y-5=0$.
}
\end{ex}

%cau 4
\begin{ex}%[0H9H3-3]
Cho hai đường thẳng $\left(d_1\right)\colon m x+y=m+1$, $\left(d_2\right)\colon x+m y=2$ song song nhau khi và chỉ khi
\choice
{$m=2$}
{$m= \pm 1$}
{$m=1$}
{\True $m=-1$}
\loigiai{
Hai đường thẳng $\left(d_1\right)\colon m x+y=m+1$, $\left(d_2\right)\colon x+m y=2$ lần lượt nhận $ \overrightarrow{n}_1=(m;1) $ và $ \overrightarrow{n}_2=(1;m) $ làm vectơ pháp tuyến.
Để hai đường thẳng này song song với nhau thì
$$
	\heva{&m^2=1 \\ &m^2+m \neq 2}
	\Leftrightarrow
	\heva{&\hoac{&m=1 \\ &m=-1} \\&\heva{&m \neq 1 \\ &m \neq-2}}
	\Leftrightarrow
	m=-1.
$$
}
\end{ex}

%cau 5
\begin{ex}%[0H9H3-3]
Với giá trị nào của $m$ thì hai đường thẳng $\left(\Delta_1\right)\colon 3 x+4 y-1=0$ và $\left(\Delta_2\right)\colon(2 m-1) x+m^2 y+1=0$ trùng nhau.
\choice
{$m=2$}
{mọi $m$}
{\True không có $m$}
{$m= \pm 1$}
\loigiai{
Để $ \Delta_1 $ trùng $ \Delta_2 $ thì
$$
	\dfrac{2m-1}{3} = \dfrac{m^2}{4} = \dfrac{1}{-1}
	\Leftrightarrow
	\heva{&-2m+1=3 \\&-m^2=4 \quad (\text{vô lý vì $ m^2 \ge 0 $})}
	\Leftrightarrow
	m \in \varnothing.
$$
}
\end{ex}

%cau 6
\begin{ex}%[0H9V3-3]
Nếu ba đường thẳng $d_1\colon 2 x+y-4=0$, $d_2\colon 5 x-2 y+3=0$ và $d_3\colon m x+3 y-2=0$ đồng quy thì $m$ nhận giá trị nào sau đây?
\choice
{$\dfrac{12}{5}$}
{$-\dfrac{12}{5}$}
{$ 12 $}
{\True $ -12 $}
\loigiai{
Xét hệ phương trình tọa độ giao điểm $ A $ của $ d_1 $ và $ d_2 $
$$
	\heva{&2 x + y - 4 = 0 \\&5 x - 2 y + 3 = 0}
	\Leftrightarrow
	\heva{&x=\dfrac{5}{9} \\&y=\dfrac{26}{9}}
	\Rightarrow
	A\left(\dfrac{5}{9} ; \dfrac{26}{9}\right).
$$
Để $ 3 $ đường thẳng đã cho đồng quy thì
$$
	A \in d_3
	\Leftrightarrow
	\dfrac{5 m}{9}+\dfrac{26}{3}-2=0 
	\Leftrightarrow
	m=-12.
$$
}
\end{ex}

%cau 7
\begin{ex}%[0H9H3-3]
Cho đường thẳng $d_1\colon 2 x+3 y+15=0$ và $d_2\colon x-2 y-3=0$. Khẳng định nào sau đây đúng?
\choice
{\True $d_1$ và $d_2$ cắt nhau và không vuông góc với nhau}
{$d_1$ và $d_2$ song song với nhau}
{$d_1$ và $d_2$ trùng nhau}
{$d_1$ và $d_2$ vuông góc với nhau}
\loigiai{
Đường thẳng $d_1v 2 x+3 y+15=0$ có một vectơ pháp tuyến là $\overrightarrow{n_1}=(2 ; 3)$ và đường thẳng $d_2\colon x-2 y-3=0$ có một vectơ pháp tuyến là $\overrightarrow{n_2}=(1 ;-2)$.
\\
Ta thấy $\dfrac{2}{1} \neq \dfrac{3}{-2}$ và $\overrightarrow{n_1} \cdot \overrightarrow{n_2}=2 \cdot 1+3 \cdot (-2)=-4 \neq 0$.
\\
Vậy $d_1$ và $d_2$ cắt nhau và không vuông góc với nhau.
}
\end{ex}

%cau 8
\begin{ex}%[0H9H3-3]
Với giá trị nào của $m$ thì hai đường thẳng $d_1\colon 3 x+4 y+10=0$ và $d_2\colon(2 m-1) x+m^2 y+10=0$ trùng nhau?
\choice
{$m \pm 2$}
{$m= \pm 1$}
{\True $m=2$}
{$m=-2$}
\loigiai{
Để $ d_1 $ và $ d_2 $ trùng nhau thì 
$$
	\dfrac{2 m-1}{3}=\dfrac{m^2}{4}=\dfrac{10}{10}
	\Leftrightarrow
	\heva{&2 m-1=3 \\&m^2=4}
	\Leftrightarrow
	m=2.
$$
}
\end{ex}

%cau 9
\begin{ex}%[0H9H3-3]
Tìm $m$ để hai đường thẳng $d_1\colon 2 x-3 y+4=0$ và $d_2\colon\heva{&x=2-3 t \\ &y=1-4 m t}$ cắt nhau.
\choice
{$m \neq-\dfrac{1}{2}$}
{$m \neq 2$}
{\True $m \neq \dfrac{1}{2}$}
{$m=\dfrac{1}{2}$}
\loigiai{
Đường thẳng $ d_1 $ có một vectơ pháp tuyến là $ \overrightarrow{n}_1=(2;-3) $.
\\
Đường thẳng $ d_2 $ có một vectơ chỉ phương là $ \overrightarrow{u}_2=(-3;-4m) $ và một vectơ pháp tuyến $ \overrightarrow{n}_2=(4m;-3) $.	
\\
Để $ d_1 $ và $ d_2 $ cắt nhau thì
$$
	2 \cdot (-3) \ne 4m \cdot (-3)
	\Leftrightarrow
	m \neq \dfrac{1}{2}.
$$
}
\end{ex}

%cau 10
\begin{ex}%[0H9H3-3]
Góc giữa hai đường thẳng $a \colon \sqrt{3} x-y+7=0$ và $b\colon x-\sqrt{3} y-1=0$ là
\choice
{\True $30^{\circ}$}
{$90^{\circ}$}
{$60^{\circ}$}
{$45^{\circ}$}
\loigiai{
Đường thẳng $a$ có vectơ pháp tuyến là $\overrightarrow{n_1}=(\sqrt{3} ;-1)$;
\\
Đường thẳng $b$ có vectơ pháp tuyến là $\overrightarrow{n_2}=(1 ;-\sqrt{3})$.
\\
Áp dụng công thức tính góc giữa hai đường thẳng, ta có
$$
	\cos (a, b)
	=\dfrac{\left|\overrightarrow{n_1} \cdot \overrightarrow{n_2}\right|}{\left|\overrightarrow{n_1}\right| \cdot \left|\overrightarrow{n_2}\right|}
	=\dfrac{\left|1 \cdot \sqrt{3}+(-1)(-\sqrt{3})\right|}{2 \cdot 2}
	=\dfrac{\sqrt{3}}{2}.
$$
Suy ra góc giữa hai đường thẳng bằng $30^{\circ}$.
}
\end{ex}

%cau 11
\begin{ex}%[0H9H3-3]
Tìm côsin góc giữa hai đường thẳng $\Delta_1\colon 2 x+y-1=0$ và $\Delta_2\colon\heva{&x=2+t \\ &y=1-t}$.
\choice
{$\dfrac{\sqrt{10}}{10}$}
{$\dfrac{3}{10}$}
{$\dfrac{3}{5}$}
{\True $\dfrac{3 \sqrt{10}}{10}$}
\loigiai{
Véctơ pháp tuyến của đường thẳng $\Delta_1$ là $\overrightarrow{n}=(2 ; 1)$ nên véctơ chỉ phương $\overrightarrow{u}=(1 ;-2)$.
\\
Véctơ chỉ phương của đường thẳng $\Delta_2$ là $\overrightarrow{u'}=(1 ;-1)$.
\\
Khi đó
$$
	\cos \left(\Delta_1 ; \Delta_2\right)
	=\left|\cos \left(\overrightarrow{u} ; \overrightarrow{u'}\right)\right|
	=\dfrac{\left|\overrightarrow{u} \cdot \overrightarrow{u'}\right|}{|\overrightarrow{u}| \cdot\left|\overrightarrow{u'}\right|}
	=\dfrac{3}{\sqrt{5} \cdot \sqrt{2}}
	=\dfrac{3 \sqrt{10}}{10}.
$$
}
\end{ex}

%cau 12
\begin{ex}%[0H9H3-3]
Cho đường thẳng $d \colon 3 x+5 y-15=0$. Trong các điểm sau đây, điểm nào \textbf{không} thuộc đường thẳng $d$?
\choice
{$M_1(5 ; 0)$}
{$M_4(-5 ; 6)$}
{$M_2(0 ; 3)$}
{\True $M_3(5 ; 3)$}
\loigiai{
Thay tọa độ các điểm vào phương trình đường thẳng $d$, ta có $M_1\,, M_4,\, M_2 \in d$ và $M_3 \notin d$.
}
\end{ex}


\Closesolutionfile{ans}
\indapan{6}{ans/ans-0-C7B10-De2-LC}

\cauds
\Opensolutionfile{ans}[ans/ans-0-C7B10-De2-DS]
%cau 1
\begin{ex}%[0H9H3-3]
Cho hai đường thẳng $ \Delta_1\colon 2 x+y+15=0 $ và $ \Delta_2\colon x-2 y-3=0 $. Các mệnh đề sau đúng hay sai?
\choiceTF
{\True $\Delta_1$ có vectơ pháp tuyến $\overrightarrow{n}_1=(2 ; 1)$, $\Delta_2$ có vectơ pháp tuyến $\overrightarrow{n}_2=(1 ;-2)$}
{\True Hai đường thẳng $\Delta_1$, $\Delta_2$ cắt nhau}
{$\Delta_1$, $\Delta_2$ cắt nhau tại $\left(-\dfrac{27}{4} ;-\dfrac{21}{4}\right)$}
{\True $\Delta_1$, $\Delta_2$ vuông góc với nhau}
\loigiai{
\begin{itemchoice}
	\itemch $\Delta_1$ có vectơ pháp tuyến $\overrightarrow{n}_1=(2 ; 1)$, $\Delta_2$ có vectơ pháp tuyến $\overrightarrow{n}_2=(1 ;-2)$.
	
	\itemch Vì $2 \cdot (-2) \neq 1 \cdot 1$ nên hai vectơ trên không cùng phương, suy ra hai đường thẳng $\Delta_1$, $\Delta_2$ cắt nhau.
	
	\itemch Xét hệ
	$$
		\heva{&2 x+y+15=0 \\ &x-2 y-3=0}
		\Leftrightarrow
		\heva{&x=-\dfrac{27}{5} \\ &y=-\dfrac{21}{5}}
	$$
	Vậy $\Delta_1$, $\Delta_2$ cắt nhau tại $\left(-\dfrac{27}{5} ;-\dfrac{21}{5}\right)$.
	
	\itemch Ta thấy
	$\overrightarrow{n}_1 \cdot \overrightarrow{n}_2=2 \cdot 1+1 \cdot(-2)=0$.
	\\
	Vậy $d_1$ và $d_2$ vuông góc với nhau.
\end{itemchoice}
}
\end{ex}

%cau 2
\begin{ex}%[0H9H3-3]
Cho đường thẳng $d\colon x+2 y-1=0$. Các mệnh đề sau đúng hay sai?
\choiceTF
{\True $d$ cắt $\Delta_1\colon-x+3 y=0$ tại $A\left(\dfrac{3}{5} ; \dfrac{1}{5}\right)$}
{\True $d \parallel \Delta_2\colon y=-\dfrac{1}{2} x+3$}
{\True $d \parallel \Delta_3\colon 3 x+6 y+3=0$}
{$d$ trùng với $\Delta_4\colon 2 x+y-1=0$}
\loigiai{
\begin{itemchoice}
	\itemch Ta có $\dfrac{1}{-1} \neq \dfrac{2}{3}$ nên $d$ cắt $\Delta_1$. Tọa độ giao điểm của $d$ và $\Delta_1$ là nghiệm của hệ phương trình
	$$
		\heva{&x + 2 y - 1 = 0 \\&- x + 3 y = 0}
		\Leftrightarrow
		\heva{&x=\dfrac{3}{5} \\&y=\dfrac{1}{5}}
	$$
	Vậy $d$ cắt $\Delta_1$ tại $A\left(\dfrac{3}{5} ; \dfrac{1}{5}\right)$
	
	\itemch $\Delta_2\colon y=-\dfrac{1}{2} x+3 \Leftrightarrow x+2 y-6=0$.
	\\
	Ta có $\dfrac{1}{1}=\dfrac{2}{2} \neq \dfrac{-1}{-6}$ nên $d \parallel \Delta_2$.
	
	\itemch Ta có $\dfrac{1}{3}=\dfrac{2}{6} \neq \dfrac{-1}{3}$ nên $d \parallel \Delta_3$.
	
	\itemch $d\colon x + 2 y - 1 = 0$ cắt $\Delta_4 \colon 2 x+y-1=0$ vì $ \dfrac{1}{1} \ne \dfrac{2}{1} $.
\end{itemchoice}
}
\end{ex}

%cau 3
\begin{ex}%[0H9H3-3]
Cho $\Delta_1\colon x-y-3=0$, $\Delta_2\colon\heva{&x=1-t \\ &y=2+2 t}$. Các mệnh đề sau đúng hay sai?
\choiceTF
{\True $\Delta_1$ có vectơ pháp tuyến $\overrightarrow{n}_1=(-1 ;-1)$}
{$\Delta_2$ có vectơ pháp tuyến $\overrightarrow{n}_2=(2 ;-1)$}
{\True Hai đường thẳng $\Delta_1$, $\Delta_2$ cắt nhau}
{$\Delta_1$, $\Delta_2$ cắt nhau tại điểm có tọa độ $\left(\dfrac{7}{2} ;-\dfrac{2}{3}\right)$}
\loigiai{
\begin{itemchoice}
	\itemch $\Delta_1\colon x-y-3=0$ có vectơ pháp tuyến $\overrightarrow{n}_1=(-1 ;-1)$.
	
	\itemch $\Delta_2\colon\heva{&x=1-t \\ &y=2+2 t}$ có vectơ chỉ phương $\overrightarrow{u}_2=(-1;2)$ và vectơ pháp tuyến $\overrightarrow{n}_2=(2 ;1)$.
	
	\itemch Xét $\overrightarrow{n}_1=(-1 ;-1)$ và $\overrightarrow{n}_2=(2 ;1)$ có $-1 \cdot 1 \neq-1 \cdot 2$ nên hai vectơ này không cùng phương. Do đó hai đường thẳng $\Delta_1$, $\Delta_2$ cắt nhau.
	
	\itemch Thay phương trình $\Delta_2$ vào phương trình $\Delta_1:(1-t)-(2+2 t)-3=0$, ta được
	$$
		-3 t-4=0
		\Leftrightarrow
		t=-\dfrac{4}{3}
		\Rightarrow
		\heva{&x=\dfrac{7}{3} \\&y=-\dfrac{2}{3}.}
	$$
	Vậy $\Delta_1$, $\Delta_2$ cắt nhau tại điểm có tọa độ $\left(\dfrac{7}{3} ;-\dfrac{2}{3}\right)$.
\end{itemchoice}
}
\end{ex}

%cau 4
\begin{ex}%[0H9V3-3]
Các mệnh đề sau đúng hay sai?
\choiceTF
{$d_1\colon x+\sqrt{3} y=0, d_2\colon x+10=0$ có $\left(d_1, d_2\right)=45^{\circ}$}
{$d_1\colon 2 x+2 \sqrt{3} y+\sqrt{5}=0$, $d_2\colon y-\sqrt{6}=0$ có $\left(d_1, d_2\right)=60^{\circ}$}
{$\Delta_1\colon\heva{&x=4+2 t \\ &y=1-3 t}$ và $\Delta_2\colon 3 x+2 y-14=0$ có $\left(\Delta_1, \Delta_2\right)=30^{\circ}$}
{\True $\Delta_1\colon x-3 y+3=0$ và $\Delta_2\colon x-3 y-5=0$ có $\Delta_1 \parallel \Delta_2$}
\loigiai{
\begin{itemchoice}
	\itemch Hai đường thẳng $d_1$, $d_2$ có cặp vectơ pháp tuyến $\overrightarrow{n}_1=(1 ; \sqrt{3})$, $\overrightarrow{n}_2=(1 ; 0)$.
	\\
	Vì vậy
	$$
		\cos \left(d_1, d_2\right)
		=\dfrac{\left|\overrightarrow{n}_1 \cdot \overrightarrow{n}_2\right|}{\left|\overrightarrow{n}_1\right| \cdot\left|\overrightarrow{n}_2\right|}
		=\dfrac{|1.1+\sqrt{3} \cdot 0|}{\sqrt{1+3} \cdot \sqrt{1+0}}
		=\dfrac{1}{2}.
	$$
	Suy ra $\left(d_1, d_2\right)=60^{\circ}$.
	
	\itemch Hai đường thẳng $d_1$, $d_2$ có cặp vectơ pháp tuyến $\overrightarrow{n}_1=(2 ; 2 \sqrt{3})$, $\overrightarrow{n}_2=(0 ; 1)$.
	\\
	Vì vậy
	$$
		\cos \left(d_1, d_2\right)
		=\dfrac{\left|\overrightarrow{n}_1 \cdot \overrightarrow{n}_2\right|}{\left|\overrightarrow{n}_1\right| \cdot\left|\overrightarrow{n}_2\right|}
		=\dfrac{|2.0+2 \sqrt{3} \cdot 1|}{\sqrt{4+12} \cdot \sqrt{0+1}}
		=\dfrac{\sqrt{3}}{2}.
	$$
	Suy ra $\left(d_1, d_2\right)=30^{\circ}$.
	
	\itemch $\Delta_1$ có vectơ chỉ phương $\overrightarrow{u}_1=(2 ;-3)$ nên có một vectơ pháp tuyến $\overrightarrow{n}_1=(3 ; 2)$;
	\\
	$\Delta_2$ có một vectơ pháp tuyến $\overrightarrow{n}_2=(3 ; 2)$.
	\\
	Ta có $3 \cdot 2=2 \cdot 3$ nên hai vectơ pháp tuyến này cùng phương nhau.
	\\
	Mặt khác, điểm $A(4 ; 1) \in d_1$ và $A \in d_2$.
	Vậy $\Delta_1$, $\Delta_2$ trùng nhau.
	
	\itemch Hai đường thẳng $\Delta_1$, $\Delta_2$ lần lượt có vectơ pháp tuyến $\overrightarrow{n}_1=(1 ;-3)$, $\overrightarrow{n}_2=(1 ;-3)$ với $1 \cdot (-3)=-3 \cdot 1$ nên hai vectơ này cùng phương.
	\\
	Mặt khác, $A(0 ; 1) \in \Delta_1$ mà $A \notin \Delta_2$ nên hai đường thẳng này song song nhau.
\end{itemchoice}
}
\end{ex}

\Closesolutionfile{ans}
\indapan{2}{ans/ans-0-C7B10-De2-DS}


\Opensolutionfile{ans}[ans/ans-0-C7B10-De2-KQ]
\caukq
%cau 1
\begin{ex}%[0H9H3-3]
Đường thẳng $d$ song song và cách đường thẳng $\Delta: y-3=0$ một khoảng bằng $ 5 $. Tìm được bao nhiêu đường thẳng $ d $ thỏa mãn yêu cầu bài toán?
\shortans{$ 2 $}
\loigiai{
Ta có $d \parallel \Delta: y-3=0 \Rightarrow$ Phương trình $d$ có dạng $y+c=0$.
\\
Ta có $M(0 ; 3) \in \Delta$. Vì $d$ cách $\Delta$ một khoảng bằng 5 nên
$$
	\mathrm{d}(d, \Delta)=5
	\Rightarrow \mathrm{d}(M, d)=5 
	\Rightarrow 
	\dfrac{|3+c|}{\sqrt{0+1}}=5
	\Rightarrow
	\hoac{&c=2 \\&c=-8.}
$$
Vậy có hai phương trình đường thẳng thỏa mãn là $y+2=0$; $y-8=0$.
}
\end{ex}

%cau 2
\begin{ex}%[0H9H3-3]
Trong mặt phẳng toạ độ $O x y$, cho điểm $I(-2 ; 4)$. Tính bán kính của đường tròn tâm $I$ tiếp xúc với đường thẳng $\Delta\colon\heva{&x=2+3 t \\ &y=-2-t}$. (Làm tròn kết quả đến hàng phân mười).
\shortans{$ 4{,}4 $}
\loigiai{
Đường thẳng $\Delta\colon\heva{&x=2+3 t \\ &y=-2-t}$ có vectơ chỉ phương là $\overrightarrow{u}=(3 ;-1)$ nên nhận $\overrightarrow{n}=(1 ; 3)$ làm vectơ pháp tuyến.
\\
Do đó, phương trình tổng quát của đường thẳng $\Delta$ là $(x-2)+3(y+2)=0 \Leftrightarrow x+3 y+4=0$.
\\
Vì đường tròn tâm $I$ tiếp xúc với đường thẳng $\Delta$ tâm $I$ nên có bán kính $ R $ bằng khoảng cách từ $I$ đến đường thẳng $\Delta$.
$$
	R=\mathrm{d}(I, \Delta)
	=\dfrac{|(-2)+3 \cdot 4+4|}{\sqrt{1^2+3^2}} \approx 4{,}4
$$
}
\end{ex}

%cau 3
\begin{ex}%[0H9H3-3]
Trong mặt phẳng tọa độ $O x y$, cho điểm $A(-2 ; 5)$. Điểm $M$ trên trục hoành sao cho đường thẳng $\Delta\colon 3 x+2 y-3=0$ cách đều hai điểm $A$ và $M$. Tổng hoành độ $ 2 $ điểm $ M $ tìm được là bao nhiêu?
\shortans{$ 2 $}
\loigiai{
Gọi $M(a ; 0)$ là điểm thuộc trục hoành. Khoảng cách từ $A, M$ đến đường thẳng $\Delta\colon 3 x+2 y-3=0$ lần lượt là $\dfrac{1}{\sqrt{13}}$, $\dfrac{|3 a-3|}{\sqrt{13}}$.
\\
Vì đường thẳng $\Delta\colon 3 x+2 y-3=0$ cách đều hai điểm $A$, $M$ nên
$$
	\dfrac{1}{\sqrt{13}}=\dfrac{|3 a-3|}{\sqrt{13}} 
	\Leftrightarrow 
	|3 a-3|=1
	\Leftrightarrow 
	\hoac{&a=\dfrac{4}{3} \\&a=\dfrac{2}{3}.}
$$
Ta tìm được $ 2 $ điểm $M\left(\dfrac{4}{3} ; 0\right)$ hoặc $M\left(\dfrac{2}{3} ; 0\right)$.
\\
Tổng hoành độ là $ \dfrac{4}{3}+\dfrac{2}{3}=2 $.
}
\end{ex}

%cau 4
\begin{ex}%[0H9H3-3]
Ta tìm được $ 2 $ tham số $m_1$, $ m_2 $ để $ \Delta_1\colon\heva{&x=8-(m+1) t \\&y=10+t} $ và $ \Delta_2\colon m x+2 y-14=0 $ song song với nhau.
Tính $ m_1^2+m_2^2 $.
\shortans{$ 5 $}
\loigiai{
$\Delta_1$, $\Delta_2$ lần lượt có vectơ pháp tuyến $\overrightarrow{n}_1=(1 ; m+1)$, $\overrightarrow{n}_2=(m ; 2)$.
\\
Điều kiện cần: $\Delta_1 \parallel \Delta_2 \Rightarrow \overrightarrow{n}_1$ cùng phương với $\overrightarrow{n}_2$
\\
$
	\Rightarrow 1 \cdot 2=(m+1) m 
	\Rightarrow m^2+m-2=0 
	\Rightarrow \hoac{&m=1 \\&m=-2.}
$
\\
Thử lại (điều kiện đủ):
\begin{itemize}
	\item Với $m=1$ thì $\Delta_1\colon\heva{&x=8-2 t \\ &y=10+t}$, $\Delta_2\colon x+2 y-14=0$ (hai đường thẳng này đã có cặp vectơ pháp tuyến cùng phương nhau).
	\\
	Vì $A(8 ; 10) \in \Delta_1$, $A \notin \Delta_2$ nên $\Delta_1 \parallel \Delta_2$.
	\\
	Do vậy $m=1$ thỏa mãn đề bài.
	
	\item Với $m=-2$ thì $\Delta_1\colon\heva{&x=8+t \\ &y=10+t}$, $\Delta_2\colon-2 x+2 y-14=0$ (hai đường thẳng này đã có cặp vectơ pháp tuyến cùng phương nhau).
	\\
	Vì $A(8 ; 10) \in \Delta_1$, $A \notin \Delta_2$ nên $\Delta_1 \parallel \Delta_2$.
	\\
	Do vậy $m=-2$ thỏa mãn đề bài.
\end{itemize}
Ta tìm được hai giá trị $m$ thỏa mãn là $m_1=1$ và $m_2=-2$.
\\
Do vậy $ m_1^2+m_2^2=5 $.
}
\end{ex}

%cau 5
\begin{ex}%[0H9H3-3]
Có bao nhiêu giá trị của tham số $m$ để hai đường thẳng $\Delta_1: 3 x+4 y-1=0$ và $\Delta_2:(2 m-1) x+m^2 y+1=0$ trùng nhau?
\shortans{$ 0 $}
\loigiai{
$\Delta_1$, $\Delta_2$ có hai vectơ pháp tuyến là $\overrightarrow{n}_1=(3 ; 4)$, $\overrightarrow{n}_2=\left(2 m-1 ; m^2\right)$.
\\
Điều kiện cần: $\Delta_1, \Delta_2$ trùng nhau suy ra hai vectơ $\overrightarrow{n}_1$, $\overrightarrow{n}_2$ cùng phương, suy ra
$$
	3 m^2=4 \cdot(2 m-1) 
	\Leftrightarrow 3 m^2-8 m+4=0 
	\Leftrightarrow \hoac{&m=2 \\&m=\dfrac{2}{3}.}
$$
\begin{itemize}
	\item Với $m=2$ thì $\Delta_2\colon 3 x+4 y+1=0$.
	\\
	Ta thấy $A(-1 ; 1) \in \Delta_1$ mà $A \notin \Delta_2$ nên $\Delta_1$, $\Delta_2$ không trùng nhau (loại $m=2$).
	
	\item Với $m=\dfrac{2}{3}$ thì $\Delta_2\colon \dfrac{1}{3} x+\dfrac{4}{9} y+1=0$.
	\\
	Ta thấy $A(-1 ; 1) \in \Delta_1$ mà $A \notin \Delta_2$ nên $\Delta_1$, $\Delta_2$ không trùng nhau (loại $m=\dfrac{2}{3}$).
\end{itemize}
Vậy không có giá trị $m$ nào thỏa mãn đề bài.
}
\end{ex}

%cau 6
\begin{ex}%[0H9H3-3]
Tìm $m$ để hai đường thẳng $\Delta_1\colon\heva{&x=8+(m+1) t \\ &y=10-t}$ và $\Delta_2\colon m x+6 y-76=0$ song song với nhau.
\shortans{$ -3 $}
\loigiai{
Hai đường thẳng $\Delta_1$, $\Delta_2$ có cặp vectơ pháp tuyến $\overrightarrow{n}_1=(1 ; m+1)$, $\overrightarrow{n}_2=(m ; 6)$.
\\
Điều kiện cần để $\Delta_1, \Delta_2$ song song nhau là $\overrightarrow{n}_1$, $\overrightarrow{n}_2$ cùng phương
$
	\Leftrightarrow 1.6=(m+1) m 
	\Leftrightarrow m^2+m-6=0 
	\Leftrightarrow \hoac{&m=2 \\&m=-3.}
$
\begin{itemize}
	\item Với $m=2$ thì $\Delta_2\colon 2 x+6 y-76=0 \Leftrightarrow x+3 y-38=0$.
	\\
	Ta có $A(8 ; 10) \in \Delta_1$, $A \in \Delta_2$ nên loại $m=2$.
	
	\item Với $m=-3$ thì $\Delta_2\colon-3 x+6 y-76=0$.
	\\
	Ta có $A(8 ; 10) \in \Delta_1$, $A \notin \Delta_2$ nên loại $m=-3$ thỏa mãn.
\end{itemize}
Vậy với $m=-3$ thì $\Delta_1, \Delta_2$ song song nhau.
}
\end{ex}

\Closesolutionfile{ans}
\indapan{6}{ans/ans-0-C7B10-De2-KQ}